\section*{\centering{\Large Facilities and Equipment}}

\section*{Computing}
%
\hrule \vspace{0.5mm}

As a member of the University of Texas (UT) system, UT Dallas faculty get
preferential access to the supercomputers at the Texas Advanced Computing Center
(TACC) \citep{tacc_22}, one of the flagship national computing centers.
Particularly relevant for this proposal are TACC's Lonestar6
\citep{lonestar6_22} and Stampede2 \citep{stampede2_22} systems. Lonestar6 is a
cluster dedicated for use by researchers at UT institutions, Texas A\&M
University, Texas Tech University, and the University of North Texas, and
Stampede2 is part of the national computing infrastructure, but UT and other
partner institutions have a separate allocation request mechanism with dedicated
guaranteed time. Allocations on both systems are virtually guaranteed for UT
faculty given a successful funding proposal. For this academic year, the PI has
an allocation of 30\,000 node hours on Lonestar6 ($\sim 4 \times 10^6$ core
hours) and 35\,000 node hours on Stampede2 ($\sim 1.7 \times 10^6$ core hours).
Standard research allocations on theses systems are 21\,000 node hours on
stampede2 and for up to 100\,000 node hours on Lonestar6.

UT Dallas maintains an additional high performance computing cluster
(Ganymede)\citep{ganymede_22}, with 4786 cores, exclusively for UT Dallas
researchers.  There are no formal usage quotas on the cluster, with users
submitting jobs to a load balancing system consisting of a number of queues with
different maximum job run times, node allocations and priorities. PI Penev and
his students already have access to the cluster and are regularly using it for
hundreds of thousands of core-hours each year.

Finally, PI Penev has a dedicated 32 core computer only for use by his group
that will contribute tens of thousands of additional core-hours each year to this
project, mostly to serve as a convenient software development and testing
system.

The available computing resources greatly exceed the requirements for this
project. Analyzing \kepler{} eclipsing binaries in a very similar manner as we
propose to do, \citep{Windemuth_et_al_19} found that in order for the MCMC
posteriors to converge, $10^5$ --- $10^6$ steps of 128 \texttt{emcee} walkers
are required. By far the slowest part of each step is generating the lightcurve
model, and within that, modeling the transits dominates over the phase curve
model. Since we already have all the steps for modeling the transits at hand,
and even a simplified phase curve model, we were able to generate $10^6$ transit
models with random system parameters in order to estimate the required computing
resources. Scaling by the expected number of systems to analze, the steps each
system requires, and accounting for the well understood efficiency of
\texttt{emcee} multi-processing, we can accurately estimate that just under
20\,000 node hours on Lonestar6 will be sufficient to carry out the entire
eclipsing binary physiycal parameters analysis. The spin analysis requires
negligible computational resources compared to the physical parameter analysis.
As a result, this entire project will require less than 10\% of a typical
reseacrh allocation on Lonestar6 over two years.
